% Appendix: progress tracking (nodes/progress.py, nodes/progress_act.py).
\section{Progress tracking against a task agenda}
\label{app:progress}

When the meeting comes with a task agenda $\mathcal{G}$, \method{} also reports the status of each agenda task.
This path is deliberately simpler than the one for action items and decisions: updates are checked by a classifier in the evidence-check node and do not enter the Verifier.

\paragraph{Extraction.}
The Progress agent makes one LLM call per segment.
It receives the whole agenda (identifier, title, owner, supporting parties, due date, status before the meeting) and selects the tasks the segment talks about itself; there is no keyword pre-filter, because tasks are often discussed without their title being said.
For each selected task it reports a status in \{\texttt{done}, \texttt{in\_progress}, \texttt{blocked}, \texttt{returned}, \texttt{deferred}, \texttt{mentioned}\}, a status turn, evidence turns, and optionally blockers, a directive and a new deadline.

\paragraph{Checking.}
A hosted multiple-choice classifier reads the status turn (windowed when long, as in \S\ref{sec:method:evidence}) and returns a distribution over statuses, including ``only mentioned or asked'' and ``not about this task''.
A separate question first decides whether the task only asks for a report to the meeting or must produce a deliverable (a document, a submission, an approved procedure); report-only tasks are judged with their own status question (reported, not yet ready, asks to postpone, asked to report again).
\begin{itemize}[nosep,leftmargin=*]
  \item \emph{Update with a status.} If the classifier confirms the status, it is kept; if it is confident in a different status, the status is corrected; otherwise the update is downgraded to \texttt{mentioned} and flagged uncertain.
  A generic question about the task's deliverable tends to score the state of a project being mentioned rather than of the task, so it is not allowed to flip between \texttt{done}, \texttt{in\_progress} and \texttt{blocked}.
  \item \emph{Update marked \texttt{mentioned}.} If every evidence turn is classified as not about the task, the update is removed (the agent linked the wrong task).
  If a turn about the task is confidently a report (done, in progress, blocked), the status is raised accordingly.
  A request to redo or postpone is raised to \texttt{returned}/\texttt{deferred} only if the commitment-act test (\S\ref{sec:method:evidence}) also classifies that turn as a directive or conclusion, because the status classifier sees only two preceding turns and cannot tell whether a member's proposal was accepted by the chair.
\end{itemize}

\paragraph{Aggregation.}
Finalize reduces the updates of each task across all segments to one row: an update whose status turn is a directive or conclusion overrides a report; between updates of the same kind, the later one wins.
The reduction never relies on speaker names, because a diarised speaker label may cover several people.
A task without any update is reported as \texttt{not\_reported}; agenda tasks are never removed.
New action items are related to the agenda as \texttt{new}, \texttt{follow\_up}, \texttt{reassign} or \texttt{reschedule}: the Action agent may tag an item with an agenda identifier (the schema gains an enumerated \texttt{task\_id} field only when an agenda exists), and the relation is derived from the item's identifier, status, assignees and deadline compared with the agenda entry.
